\chapter{Dataset Description and Preprocessing}
\label{ch:dataset}

\section{HAM10000 Dataset Overview}

The HAM10000 (``Human Against Machine with 10,000 training images'') dataset is a large collection of multi-source dermatoscopic images of skin lesions. This dataset has become the de facto standard benchmark for evaluating skin lesion classification algorithms.

\subsection{Data Source and Collection}

The HAM10000 dataset aggregates dermatoscopic images from multiple sources:

\begin{itemize}
    \item \textbf{UNBC (University of British Columbia)}: Clinical images from dermatology clinics
    \item \textbf{BCN (Barcelona-Ussai)}: Dermoscopy dataset from Barcelona
    \item \textbf{ViDIR (Vienna Dynamic Interactive Dermoscopy Reference)}: Interactive dermoscopy reference images
    \item \textbf{Private Dermatology Clinics}: Clinical cases from practicing dermatologists
\end{itemize}

Multi-source aggregation introduces heterogeneity in imaging protocols, camera equipment, and dermatoscopic techniques, increasing dataset complexity and real-world relevance.

\section{Dataset Composition and Statistics}

\subsection{Overall Dataset Size}

The HAM10000 dataset contains a total of \textbf{10,015 dermatoscopic images} categorized into seven distinct skin lesion types:

\begin{table}[H]
\centering
\caption{HAM10000 Dataset Class Distribution}
\label{tab:dataset_distribution}
\begin{tabular}{|l|c|c|c|}
\hline
\textbf{Lesion Type} & \textbf{Code} & \textbf{Count} & \textbf{\% of Total} \\
\hline
Actinic Keratosis / Solar Keratosis & akiec & 327 & 3.27\% \\
Basal Cell Carcinoma & bcc & 514 & 5.14\% \\
Benign Keratosis-like Lesions & bkl & 1099 & 10.98\% \\
Dermatofibroma & df & 115 & 1.15\% \\
Melanoma & mel & 1113 & 11.12\% \\
Melanocytic Nevi & nv & 6705 & 67.00\% \\
Vascular Lesions & vasc & 142 & 1.42\% \\
\hline
\textbf{Total} & & \textbf{10,015} & \textbf{100\%} \\
\hline
\end{tabular}
\end{table}

\subsection{Class Imbalance Analysis}

The dataset exhibits significant class imbalance, with Melanocytic Nevi (nv) comprising 67\% of all images while Dermatofibroma represents only 1.15\%. This imbalance presents training challenges:

\begin{itemize}
    \item Models trained with standard loss functions are biased toward majority classes
    \item Minority classes (akiec, df, vasc) with fewer than 400 samples per class require special handling
    \item Threshold optimization becomes critical for achieving balanced sensitivity and specificity
\end{itemize}

The ratio between the largest class (nv: 6,705) and smallest class (df: 115) is approximately 58:1, necessitating sophisticated rebalancing strategies.

\section{Image Specifications}

\subsection{Image Properties}

Each image in the HAM10000 dataset has the following specifications:

\begin{itemize}
    \item \textbf{Format}: JPEG
    \item \textbf{Resolution}: Variable, typically 450 × 450 to 1200 × 1200 pixels
    \item \textbf{Color Space}: RGB (8-bit per channel)
    \item \textbf{Metadata}: Associated with a unique patch ID and classification code
\end{itemize}

The variable resolution requires standardization for batch processing in deep learning models.

\subsection{Image Quality Characteristics}

Dermatoscopic images exhibit specific characteristics compared to natural images:

\begin{itemize}
    \item \textbf{Close-up Detail}: Magnified views of skin surface highlighting fine morphological features
    \item \textbf{Variable Imaging Conditions}: Different dermoscope types and polarization settings create appearance variations
    \item \textbf{Artifacts}: Ruler measurements, reflections, or other non-lesion elements may be present
    \item \textbf{Illumination Variation}: Inconsistent lighting conditions across images
\end{itemize}

\section{Data Preprocessing}

\subsection{Image Resizing and Normalization}

All images are resized to a consistent $224 \times 224$ pixel resolution to:

\begin{enumerate}
    \item Enable uniform batch processing
    \item Match the input resolution of pre-trained Swin Transformer models
    \item Balance information retention and computational efficiency
\end{enumerate}

Bilinear interpolation is employed for resizing to preserve image quality during downsampling.

Image normalization uses ImageNet statistics:

\begin{equation}
\tilde{I} = \frac{I - \mu}{\sigma}
\end{equation}

where $\mu = [0.485, 0.456, 0.406]$ and $\sigma = [0.229, 0.224, 0.225]$ represent ImageNet mean and standard deviation per channel.

\subsection{Data Type Conversion}

Images are converted from integer representation (0-255) to floating-point values (0.0-1.0) before normalization to improve numerical stability during training.

\section{Data Augmentation Strategy}

Data augmentation is crucial for improving model robustness and generalization on the limited HAM10000 dataset. Our augmentation pipeline includes:

\subsection{Geometric Transformations}

\begin{itemize}
    \item \textbf{Random Rotation}: Rotations up to 180 degrees account for dermoscopic probe orientation variations
    \item \textbf{Horizontal and Vertical Flipping}: Reflect bilateral symmetries in skin lesions
    \item \textbf{Random Affine Transformations}: Translation and shearing simulate camera position variations
    \item \textbf{Elastic Deformations}: Non-linear deformations model skin surface variations
\end{itemize}

\subsection{Intensity Adjustments}

\begin{itemize}
    \item \textbf{Brightness Adjustment}: $\pm 20\%$ variation addresses illumination inconsistencies
    \item \textbf{Contrast Adjustment}: $\pm 20\%$ variation improves feature robustness
    \item \textbf{Color Jittering}: Saturation and hue adjustments simulate device calibration variations
    \item \textbf{Gaussian Noise}: Small amount ($\sigma = 0.01$) simulates sensor noise
\end{itemize}

\subsection{Advanced Augmentation Techniques}

\begin{itemize}
    \item \textbf{Mixup}: Blends pairs of images and linearly interpolates labels, preventing overfitting to specific image features
    \item \textbf{CutMix}: Cuts and pastes patches between images while interpolating labels
    \item \textbf{AutoAugment}: Learned composition of augmentation operations optimized for the specific dataset
\end{itemize}

\section{Data Splitting Strategy}

\subsection{Train-Validation-Test Split}

The dataset is partitioned using stratified random sampling to maintain class distribution:

\begin{itemize}
    \item \textbf{Training Set}: 70\% (7,010 images) - used for model parameter optimization
    \item \textbf{Validation Set}: 15\% (1,502 images) - used for hyperparameter tuning and early stopping
    \item \textbf{Test Set}: 15\% (1,503 images) - reserved for final model evaluation
\end{itemize}

Stratified sampling ensures each split maintains the original class distribution, critical for evaluating performance on imbalanced data.

\subsection{Data Leakage Prevention}

Augmentation is applied only to the training set. Validation and test sets use only center-crop and normalization to prevent data leakage of training augmentation statistics.

\section{Class Imbalance Handling}

\subsection{Random Oversampling}

During training, minority classes are randomly oversampled to balance class representation:

\begin{equation}
n_{\text{balanced}} = \max(n_1, n_2, \ldots, n_7)
\end{equation}

Each class is resampled to approximately 6,705 samples (the maximum class frequency), creating a balanced training distribution while respecting original test set distribution.

\subsection{Weighted Loss Function}

Class weights in the loss function are inversely proportional to class frequencies:

\begin{equation}
w_i = \frac{N}{7 \cdot n_i}
\end{equation}

where $N = 10,015$ is total dataset size and $n_i$ is the count of class $i$. This weighting scheme emphasizes learning from underrepresented classes.

\subsection{Focal Loss Consideration}

For highly imbalanced scenarios, focal loss is considered as an alternative:

\begin{equation}
FL(p_t) = -\alpha_t(1-p_t)^{\gamma}\log(p_t)
\end{equation}

where $\gamma = 2$ focuses on hard exemplars and $\alpha$ balances class importance.

\section{Data Quality Assurance}

\subsection{Image Inspection}

A sample of images from each class is manually inspected to verify:

\begin{itemize}
    \item Correct class assignment
    \item Adequate image quality and artifact absence
    \item Representative morphological characteristics
\end{itemize}

\subsection{Statistical Summary}

Summary statistics for image properties across the dataset:

\begin{table}[H]
\centering
\caption{Image Statistics Summary}
\label{tab:image_stats}
\begin{tabular}{|l|c|}
\hline
\textbf{Property} & \textbf{Value} \\
\hline
Mean Image Size & 756 × 756 pixels \\
Median Image Size & 720 × 720 pixels \\
Image Size Range & 450-1200 pixels \\
Mean Pixel Intensity (R) & 147.3 \\
Mean Pixel Intensity (G) & 113.8 \\
Mean Pixel Intensity (B) & 98.5 \\
\hline
\end{tabular}
\end{table}

The dataset's large size, multi-source composition, and publicly available nature make it an ideal benchmark for developing and evaluating skin lesion classification algorithms.
